% Ad Familiares 6.4 --- headnote
% ad-familiares-06-04, before mid-January 45 BC, Rome

\textit{Cicero to A. Manlius Torquatus, written at Rome before
the middle of January 45 BC (Perseus: \textit{Romae ante med.~m.
Ian.}). It is the longest and most personally pitched of the
Torquatus condolence cluster: the Spanish war is still open,
the engagement at Munda is some weeks off, and the prospect
under consideration in \S 1 --- that, whichever side wins, the
victory will not look very different to the defeated --- is
the prospect from which everything else in the letter is
written. Torquatus is in Athens, still under uncertain
clemency, and the boys in \S 3 are the children Cicero
elsewhere acknowledges as the proper object of his last efforts
on his friend's behalf (\textit{Fam.}~6.1.7).}

\textit{The philosophical centre is \S 2: that the consciousness
of upright intention --- \textit{conscientia rectae voluntatis}
--- is the greatest consolation for adverse circumstances, and
that there is no great evil except guilt. The doctrine, Stoic
in form, is one Cicero leans on across the consolatory writing
of these months; the formula here is among the most compressed
in the corpus. The personal disclosure of \S 4 is striking:
Cicero, speaking now in his own voice rather than as consoler,
brings to bear the same triple resource he had offered Torquatus
in \textit{Fam.}~6.1 --- to be without fault is to be without
the worst evil; a man whose life is nearly complete need not
fear what nature is anyway leading him to; and given who has
already fallen in this war, refusing the same fortune, if the
case demanded it, would be shameless. The pivot to the consoler
becoming the consoled is unannounced and characteristic of the
cluster.}

\textit{The valedictory in \S 5 notes that Servius Sulpicius
Rufus, on whose company Cicero had counted for Torquatus in
\textit{Fam.}~6.1.6, has now left Athens (this is the spring
that will end with Servius writing to Cicero on the death of
Tullia); the closing concession --- that Cicero will imitate
Torquatus's goodwill but not match his services --- is the
formula of an old friendship under the strain of unequal
fortunes.}
